\section{记忆在 Agent Loop 中的流转}

每当 Agent 处理一个请求时，记忆系统贯穿整个流程。核心机制是：记忆操作\textbf{环绕} LLM 调用——调用前检索，调用后写入。

典型的 Agent Loop 流程如下：

\begin{enumerate}[leftmargin=2em]
    \item 用户发送请求
    \item \textbf{Memory Retrieval}：从 External Memory 和 Episodic Memory 中检索相关记忆
    \item 将检索结果注入 Context Window（In-context Memory）
    \item LLM 基于完整上下文进行推理并生成响应
    \item \textbf{Memory Write}：将本次交互中的关键信息写入 External Memory 和 Episodic Memory
    \item 返回响应给用户
\end{enumerate}

\begin{importantbox}{模型本身是无状态的}
模型本身是无状态的；记忆系统才是制造出"有状态、有意识的 Agent"这一\textbf{幻象}的关键。
\end{importantbox}

\subsection{本章小结}

Agent Loop 中记忆操作在 LLM 调用前后进行——先检索、后写入。模型只是中间的推理引擎，记忆系统赋予它状态感。

